\section{Regression-only arms (S0--S11)}

Eleven single-head arms varying loss, normalisation, teacher-forcing
curriculum, and field representation. Headline: every arm loses to
persistence (CSI 0.043). Best regression CSI = 0.0176 (S11) ---
41\% of persistence.

\subsection{S0 --- minimal baseline}
\textbf{Loss:} L1 only. \textbf{Aug:} D4. \textbf{Norm:} BatchNorm.
\textbf{Test CSI:} 0 across all epochs; best val 0.039 at ep4 then
overfit.\\
\textbf{Finding:} The 2-layer ConvLSTM with plain L1 collapses to the
zero-field conditional median. Pred amplitude never crosses the
extreme threshold so TP=0. Confirmed the wall is structural, not a
capacity issue --- the deep 6-layer model also gets CSI \(\sim\)0.01.
\\[2pt]\noindent\includegraphics[width=\linewidth]{s0_harp11930_fullframe}\\[2pt]
\noindent\includegraphics[width=\linewidth]{S0_loss}

\subsection{S2 --- residual decode}
\textbf{Delta:} predict \(\Delta w_t\), output \(w_{t-1}+\Delta\).
\textbf{Finding:} Without teacher forcing this \emph{explodes} ---
variance ratio climbs 18\(\to\)764 across 15 epochs (autoregressive
positive feedback). Confirmed residual anchor needs a stabiliser.
(No figure: model output unusable.)

\subsection{S4 --- S2 + teacher forcing}
\textbf{Delta:} TF \(\tau{=}1{\to}0\) over training.
\textbf{Test:} CSI 0.099, HSS 0.178, SSIM 0.438 (old random split).
\textbf{Finding:} TF cures the explosion; variance ratio settles near
1. Looks ``sharp'' visually but on the fixed test set later this
sharpness turns out to be uncontrolled BatchNorm variance, not skill
(see S9). \\[2pt]\noindent\includegraphics[width=\linewidth]{s4_harp11930_fullframe}

\subsection{S5 --- S4 + 90/5/5 split + GroupNorm}
\textbf{Delta:} more training data, swap BN\(\to\)GN.
\textbf{Test:} SSIM 0.70, persistence skill +30\%, CSI \(\approx\) 0.
\textbf{Finding:} GroupNorm + more data improves spatial similarity
and beats persistence MAE, but \emph{smoother} field $\Rightarrow$
\emph{lower} CSI on extremes. Trade-off explicit: L1 rewards smoothing,
extreme classification punishes it.\\[2pt]\noindent\includegraphics[width=\linewidth]{s5_harp11930_fullframe}\\[2pt]
\noindent\includegraphics[width=\linewidth]{S5_loss}

\subsection{S3 --- S5 + 6-term composite loss}
\textbf{Loss:} L1 + SSIM + AsymExtreme + TempDiff + TempVar + extreme-pixel weight 25.
\textbf{Test:} best val CSI 0.0236 (ep13); test 0 (random split put a
flareless cube in test). \textbf{Finding:} Composite is the best
\emph{regression} on the qualitative HARP 11930 spatial integral (see
\S\ref{sec:cmp_intflux}) but the lift over single-term arms is small
($\sim$2\(\times\) over S0). S6 below isolates which term carries the
signal.\\[2pt]\noindent\includegraphics[width=\linewidth]{s3_harp11930_fullframe}\\[2pt]
\noindent\includegraphics[width=\linewidth]{S3_loss}

\subsection{S6 --- only the extreme-pixel weight term}
\textbf{Loss:} L1 + \texttt{extreme\_pixel\_weight}{=}25 (everything else off).
\textbf{Test:} val CSI 0.0215 \(\approx\) S3.
\textbf{Finding:} S3's full composite is \emph{not} additive --- 95\%
of its lift comes from the single extreme-pixel-weight term. SSIM /
asymmetric / temporal terms add \(\sim\)0. Identifies the only loss
lever that moves the needle.\\[2pt]\noindent\includegraphics[width=\linewidth]{s6_harp11930_fullframe}\\[2pt]
\noindent\includegraphics[width=\linewidth]{S6_loss}

\subsection{S8 --- S5 (L1) on fixed test cubes}
\textbf{Delta:} pin test set to \{245, 274, 49\}.
\textbf{Test CSI:} 0.0033 vs persistence 0.043 (\(13\times\) below).
\textbf{Finding:} Establishes the persistence wall on a flare-rich
fixed test set; the optimism in S4's old random-split CSI 0.099 was
a flareless test-cube artefact.\\[2pt]\noindent\includegraphics[width=\linewidth]{s8_harp11930_fullframe}\\[2pt]
\noindent\includegraphics[width=\linewidth]{S8_loss}

\subsection{S7 --- S8 + signed-asinh normalisation}
\textbf{Delta:} swap z-score for
\(\operatorname{sign}(w)\operatorname{asinh}(|w|/s)\) representation
(\textbf{not} the loss; only how data enters the model). \textbf{Test
CSI:} 0.0042 \(\approx\) S8. \textbf{Finding:} Heavy-tail input
compression has no effect on CSI. Apparent SSIM 0.98 is an
asinh-space artefact, not improvement in physical space. Representation
lever \textbf{rejected}.
\\[2pt]\noindent\includegraphics[width=\linewidth]{s7_harp11930_fullframe}\\[2pt]
\noindent\includegraphics[width=\linewidth]{S7_loss}

\subsection{S9 --- S8 with BatchNorm (= S4 on fixed test)}
\textbf{Delta:} swap GN\(\to\)BN.
\textbf{Test CSI:} 0.0118 (\(3.5\times\) S8 \emph{but} still below
persistence by 4\(\times\)). \textbf{Finding:} BN drives over-variance
(ratio 3.7, val hit 19); unstable (SSIM 0.51\(\to\)0.09, early-stop
ep5). S4's visual ``sharpness'' was uncontrolled variance, \textbf{not}
forecast skill. BN rejected.\\[2pt]\noindent\includegraphics[width=\linewidth]{s9_harp11930_fullframe}\\[2pt]
\noindent\includegraphics[width=\linewidth]{S9_loss}

\subsection{S10 --- S8 + fast TF curriculum}
\textbf{Delta:} \texttt{tf\_decay\_epochs}{=}5, patience 8.
\textbf{Test CSI:} 0.0027 \(\approx\) S8.
\textbf{Finding:} Free-running TF curriculum bumps variance ratio to
0.75 briefly at ep7 but L1 drags it back to 0.028 by ep15. TF is
necessary-but-insufficient on its own; L1 mean-collapse defeats the
de-smoothing.\\[2pt]\noindent\includegraphics[width=\linewidth]{s10_harp11930_fullframe}\\[2pt]
\noindent\includegraphics[width=\linewidth]{S10_loss}

\subsection{S11 --- S10 fast-TF \(+\) extreme weight (both working levers)}
\textbf{Delta:} S10 + S6's extreme\_pixel\_weight 25.
\textbf{Test CSI:} \textbf{0.0176} --- best of the S0--S11 series
(5\(\times\) S8, 41\% of persistence). \textbf{Finding:} The two
working levers are additive in CSI but the combo \emph{still} loses to
persistence by 2.4\(\times\). Variance ratio re-collapses to 0.036 by
ep15 --- L1 wins long-run.\\[2pt]\noindent\includegraphics[width=\linewidth]{s11_harp11930_fullframe}\\[2pt]
\noindent\includegraphics[width=\linewidth]{S11_loss}
\label{fig:s11_loss}

\subsection*{Verdict for the regression-only family}
Four levers explored (loss reweighting, BN/GN, TF curriculum,
input representation). The only loss term that moves CSI is
\texttt{extreme\_pixel\_weight}; the only TF schedule that helps is
the full \(\tau{:}1\!\to\!0\). Stacking both (S11) gets to 41\% of
persistence. \textbf{All single-head paths exhausted}; the model
cannot push regression amplitude past the extreme threshold without
also losing the L1 fit elsewhere.
\textbf{Conclusion: structural pivot needed.}
